\chapter{Duality, Riemann--Roch, and arithmetic}
\label{par-chap-rr-arithmetic}

\begin{theorem}[Adelic duality]
\label{par-thm-duality}
\dcref{1.26}
\uses{par-def-L-space,par-def-differential-divisor,par-thm-canonical-scaling,par-thm-differentials-one-dimensional}
\source{papaioannou-algebraic-rr:page-0013}
Let $D\in\Div(\mathbb F)$ and let $W=(\omega)$ be canonical. The map
\[
\mu:L(W-D)\longrightarrow\Omega_{\mathbb F}(D),\qquad x\longmapsto x\omega,
\]
is an isomorphism of $\mathbb K$-vector spaces. In particular,
$i(D)=\dim(W-D)$.
\end{theorem}
\begin{proof}
For $x\in L(W-D)$, Theorem~1.25 gives
$(x\omega)=(x)+W\geq-(W-D)+W=D$, so $\mu$ is well-defined, linear, and
injective. If $\omega_1\in\Omega(D)$ (the source prints $\omega_{\mathbb F}(D)$),
one-dimensionality gives
$\omega_1=x\omega$; then $(x)+W=(\omega_1)\geq D$, so
$x\in L(W-D)$. Thus $\mu$ is surjective. Taking dimensions and using
Lemma~1.22 gives $i(D)=\dim(W-D)$.
\end{proof}

\begin{theorem}[Riemann--Roch]
\label{par-thm-riemann-roch}
\dcref{1.27}
\uses{par-def-genus,par-def-L-space,par-thm-adelic-preliminary,par-thm-duality}
\source{papaioannou-algebraic-rr:page-0013}
If $W$ is a canonical divisor, then for every $D\in\Div(\mathbb F)$,
\[
\dim D=\deg D+1-g+\dim(W-D).
\]
\end{theorem}
\begin{proof}
Substitute $i(D)=\dim(W-D)$ from Theorem~1.26 into the adelic preliminary
formula of Theorem~1.18.
\end{proof}

\begin{corollary}[Degree and dimension of a canonical divisor]
\label{par-cor-canonical-degree}
\dcref{1.28}
\uses{par-thm-riemann-roch,par-def-differential-divisor}
\source{papaioannou-algebraic-rr:page-0013}
Every canonical divisor $W$ has $\deg W=2g-2$ and $\dim W=g$.
\end{corollary}
\begin{proof}
Set $D=0$ in Riemann--Roch to obtain $\dim W=g$; then set $D=W$ to obtain
$g=\deg W+1-g+\dim0=\deg W+1-g+1$, hence $\deg W=2g-2$.
\end{proof}

\begin{corollary}[Sharp high-degree formula]
\label{par-cor-high-degree-rr}
\dcref{1.29}
\uses{par-thm-riemann-roch,par-cor-canonical-degree,par-cor-equivalent-divisors}
\source{papaioannou-algebraic-rr:page-0013}
If $\deg D\geq2g-1$, then
\[
\dim D=\deg D+1-g.
\]
\end{corollary}
\begin{proof}
For canonical $W$, $\deg(W-D)<0$ by Corollary~1.28, so
$\dim(W-D)=0$ by Corollary~1.16. Apply Riemann--Roch.
\end{proof}

\section{An arithmetic theory for function fields}

\begin{conjecture}[ABC I]
\label{par-conj-abc-I}
\dcref{2.1}
\source{papaioannou-algebraic-rr:page-0013}
If pairwise prime integers $A,B,C$ satisfy $A+B=C$, then for every $\epsilon>0$
there is $M_\epsilon\in\mathbb N$ such that
\[
\max\{|A|,|B|,|C|\}\leq
\left(\prod_{p\mid ABC}p\right)^{1+\epsilon},
\]
where $p$ is positive and prime.
\end{conjecture}

\begin{definition}[Height and separable degree]
\label{par-def-height-separable-degree}
\source{papaioannou-algebraic-rr:page-0014}
For $r=m/n\in\mathbb Q$ in lowest terms,
$\operatorname{ht}(r)=\max\{\log|m|,\log|n|\}$. For a nonconstant
$u\in\mathbb F^*$, let $M$ be the maximal separable extension of $\mathbb K(u)$
inside $\mathbb F$; its degree $[M:\mathbb K(u)]$ is the separable degree of
$u$, denoted $\operatorname{ord}_s(u)$.
\end{definition}

\begin{conjecture}[ABC II]
\label{par-conj-abc-II}
\dcref{2.2}
\uses{par-def-height-separable-degree}
\source{papaioannou-algebraic-rr:page-0014}
For every $\epsilon>0$ there is $m_\epsilon$ such that rational $u,v$ with
$u+v=1$ satisfy
\[
\max\{\operatorname{ht}(v),\operatorname{ht}(u)\}
\leq m_\epsilon+(1+\epsilon)\sum_{p\mid ABC}\log p.
\]
The sum is printed with $ABC$ although the variables in this reformulation are
$u,v$; that visible mismatch is retained.
\end{conjecture}

\begin{theorem}[Function-field ABC]
\label{par-thm-function-field-abc}
\dcref{2.1}
\uses{par-def-genus,par-def-divisor,par-def-height-separable-degree}
\source{papaioannou-algebraic-rr:page-0014}
For a function field $\mathbb F/\mathbb K$ and $u,v\in\mathbb F^*$ with
$u+v=1$,
\[
\operatorname{deg}_s(v)=\operatorname{deg}_s(u)\leq
2g_{\mathbb K}+\sum_{P\in\Supp(A+B+C)}\deg_{\mathbb K}P,
\]
where $A,B$ are the zero divisors of $v,u$ and $C$ is their common pole divisor.
\end{theorem}
\begin{proof}
\notready
The source states this function-field ABC theorem without proof and notes that
it implies the function-field analogues of classical consequences, including
Fermat's Last Theorem.
\end{proof}

\begin{theorem}[Hurwitz formula]
\label{par-thm-hurwitz}
\dcref{2.2}
\uses{par-def-genus,par-def-divisor}
\source{papaioannou-algebraic-rr:page-0014}
For a finite separable extension of function fields $\mathbb L/\mathbb K$,
\[
2g_{\mathbb L}-2=[\mathbb L:\mathbb K](2g_{\mathbb K}-2)
 +\deg_{\mathbb L}D_{\mathbb L/\mathbb K}.
\]
More generally,
\[
2g_{\mathbb L}-2\geq[\mathbb L:\mathbb K](2g_{\mathbb K}-2)
 +\sum_{\mathfrak p}(e(\mathfrak p/P)-1)\deg_{\mathbb L}\mathfrak p,
\]
where the sum ranges over primes $\mathfrak p$ over $\mathbb L$; equality holds
if and only if every ramified prime is tamely ramified.
\end{theorem}
\begin{proof}
\notready
The source presents Hurwitz as an important input to the proof of ABC and gives
no proof here. Undefined arithmetic terms have their standard function-field
meanings.
\end{proof}
